\documentclass{article} % ICLR 2027 submission
\usepackage{iclr2027_conference,times}

%%%%% NEW MATH DEFINITIONS %%%%%

\usepackage{amsmath,amsfonts,bm}

% Mark sections of captions for referring to divisions of figures
\newcommand{\figleft}{{\em (Left)}}
\newcommand{\figcenter}{{\em (Center)}}
\newcommand{\figright}{{\em (Right)}}
\newcommand{\figtop}{{\em (Top)}}
\newcommand{\figbottom}{{\em (Bottom)}}
\newcommand{\captiona}{{\em (a)}}
\newcommand{\captionb}{{\em (b)}}
\newcommand{\captionc}{{\em (c)}}
\newcommand{\captiond}{{\em (d)}}

% Highlight a newly defined term
\newcommand{\newterm}[1]{{\bf #1}}


% Figure reference, lower-case.
\def\figref#1{figure~\ref{#1}}
% Figure reference, capital. For start of sentence
\def\Figref#1{Figure~\ref{#1}}
\def\twofigref#1#2{figures \ref{#1} and \ref{#2}}
\def\quadfigref#1#2#3#4{figures \ref{#1}, \ref{#2}, \ref{#3} and \ref{#4}}
% Section reference, lower-case.
\def\secref#1{section~\ref{#1}}
% Section reference, capital.
\def\Secref#1{Section~\ref{#1}}
% Reference to two sections.
\def\twosecrefs#1#2{sections \ref{#1} and \ref{#2}}
% Reference to three sections.
\def\secrefs#1#2#3{sections \ref{#1}, \ref{#2} and \ref{#3}}
% Reference to an equation, lower-case.
\def\eqref#1{equation~\ref{#1}}
% Reference to an equation, upper case
\def\Eqref#1{Equation~\ref{#1}}
% A raw reference to an equation---avoid using if possible
\def\plaineqref#1{\ref{#1}}
% Reference to a chapter, lower-case.
\def\chapref#1{chapter~\ref{#1}}
% Reference to an equation, upper case.
\def\Chapref#1{Chapter~\ref{#1}}
% Reference to a range of chapters
\def\rangechapref#1#2{chapters\ref{#1}--\ref{#2}}
% Reference to an algorithm, lower-case.
\def\algref#1{algorithm~\ref{#1}}
% Reference to an algorithm, upper case.
\def\Algref#1{Algorithm~\ref{#1}}
\def\twoalgref#1#2{algorithms \ref{#1} and \ref{#2}}
\def\Twoalgref#1#2{Algorithms \ref{#1} and \ref{#2}}
% Reference to a part, lower case
\def\partref#1{part~\ref{#1}}
% Reference to a part, upper case
\def\Partref#1{Part~\ref{#1}}
\def\twopartref#1#2{parts \ref{#1} and \ref{#2}}

\def\ceil#1{\lceil #1 \rceil}
\def\floor#1{\lfloor #1 \rfloor}
\def\1{\bm{1}}
\newcommand{\train}{\mathcal{D}}
\newcommand{\valid}{\mathcal{D_{\mathrm{valid}}}}
\newcommand{\test}{\mathcal{D_{\mathrm{test}}}}

\def\eps{{\epsilon}}


% Random variables
\def\reta{{\textnormal{$\eta$}}}
\def\ra{{\textnormal{a}}}
\def\rb{{\textnormal{b}}}
\def\rc{{\textnormal{c}}}
\def\rd{{\textnormal{d}}}
\def\re{{\textnormal{e}}}
\def\rf{{\textnormal{f}}}
\def\rg{{\textnormal{g}}}
\def\rh{{\textnormal{h}}}
\def\ri{{\textnormal{i}}}
\def\rj{{\textnormal{j}}}
\def\rk{{\textnormal{k}}}
\def\rl{{\textnormal{l}}}
% rm is already a command, just don't name any random variables m
\def\rn{{\textnormal{n}}}
\def\ro{{\textnormal{o}}}
\def\rp{{\textnormal{p}}}
\def\rq{{\textnormal{q}}}
\def\rr{{\textnormal{r}}}
\def\rs{{\textnormal{s}}}
\def\rt{{\textnormal{t}}}
\def\ru{{\textnormal{u}}}
\def\rv{{\textnormal{v}}}
\def\rw{{\textnormal{w}}}
\def\rx{{\textnormal{x}}}
\def\ry{{\textnormal{y}}}
\def\rz{{\textnormal{z}}}

% Random vectors
\def\rvepsilon{{\mathbf{\epsilon}}}
\def\rvtheta{{\mathbf{\theta}}}
\def\rva{{\mathbf{a}}}
\def\rvb{{\mathbf{b}}}
\def\rvc{{\mathbf{c}}}
\def\rvd{{\mathbf{d}}}
\def\rve{{\mathbf{e}}}
\def\rvf{{\mathbf{f}}}
\def\rvg{{\mathbf{g}}}
\def\rvh{{\mathbf{h}}}
\def\rvu{{\mathbf{i}}}
\def\rvj{{\mathbf{j}}}
\def\rvk{{\mathbf{k}}}
\def\rvl{{\mathbf{l}}}
\def\rvm{{\mathbf{m}}}
\def\rvn{{\mathbf{n}}}
\def\rvo{{\mathbf{o}}}
\def\rvp{{\mathbf{p}}}
\def\rvq{{\mathbf{q}}}
\def\rvr{{\mathbf{r}}}
\def\rvs{{\mathbf{s}}}
\def\rvt{{\mathbf{t}}}
\def\rvu{{\mathbf{u}}}
\def\rvv{{\mathbf{v}}}
\def\rvw{{\mathbf{w}}}
\def\rvx{{\mathbf{x}}}
\def\rvy{{\mathbf{y}}}
\def\rvz{{\mathbf{z}}}

% Elements of random vectors
\def\erva{{\textnormal{a}}}
\def\ervb{{\textnormal{b}}}
\def\ervc{{\textnormal{c}}}
\def\ervd{{\textnormal{d}}}
\def\erve{{\textnormal{e}}}
\def\ervf{{\textnormal{f}}}
\def\ervg{{\textnormal{g}}}
\def\ervh{{\textnormal{h}}}
\def\ervi{{\textnormal{i}}}
\def\ervj{{\textnormal{j}}}
\def\ervk{{\textnormal{k}}}
\def\ervl{{\textnormal{l}}}
\def\ervm{{\textnormal{m}}}
\def\ervn{{\textnormal{n}}}
\def\ervo{{\textnormal{o}}}
\def\ervp{{\textnormal{p}}}
\def\ervq{{\textnormal{q}}}
\def\ervr{{\textnormal{r}}}
\def\ervs{{\textnormal{s}}}
\def\ervt{{\textnormal{t}}}
\def\ervu{{\textnormal{u}}}
\def\ervv{{\textnormal{v}}}
\def\ervw{{\textnormal{w}}}
\def\ervx{{\textnormal{x}}}
\def\ervy{{\textnormal{y}}}
\def\ervz{{\textnormal{z}}}

% Random matrices
\def\rmA{{\mathbf{A}}}
\def\rmB{{\mathbf{B}}}
\def\rmC{{\mathbf{C}}}
\def\rmD{{\mathbf{D}}}
\def\rmE{{\mathbf{E}}}
\def\rmF{{\mathbf{F}}}
\def\rmG{{\mathbf{G}}}
\def\rmH{{\mathbf{H}}}
\def\rmI{{\mathbf{I}}}
\def\rmJ{{\mathbf{J}}}
\def\rmK{{\mathbf{K}}}
\def\rmL{{\mathbf{L}}}
\def\rmM{{\mathbf{M}}}
\def\rmN{{\mathbf{N}}}
\def\rmO{{\mathbf{O}}}
\def\rmP{{\mathbf{P}}}
\def\rmQ{{\mathbf{Q}}}
\def\rmR{{\mathbf{R}}}
\def\rmS{{\mathbf{S}}}
\def\rmT{{\mathbf{T}}}
\def\rmU{{\mathbf{U}}}
\def\rmV{{\mathbf{V}}}
\def\rmW{{\mathbf{W}}}
\def\rmX{{\mathbf{X}}}
\def\rmY{{\mathbf{Y}}}
\def\rmZ{{\mathbf{Z}}}

% Elements of random matrices
\def\ermA{{\textnormal{A}}}
\def\ermB{{\textnormal{B}}}
\def\ermC{{\textnormal{C}}}
\def\ermD{{\textnormal{D}}}
\def\ermE{{\textnormal{E}}}
\def\ermF{{\textnormal{F}}}
\def\ermG{{\textnormal{G}}}
\def\ermH{{\textnormal{H}}}
\def\ermI{{\textnormal{I}}}
\def\ermJ{{\textnormal{J}}}
\def\ermK{{\textnormal{K}}}
\def\ermL{{\textnormal{L}}}
\def\ermM{{\textnormal{M}}}
\def\ermN{{\textnormal{N}}}
\def\ermO{{\textnormal{O}}}
\def\ermP{{\textnormal{P}}}
\def\ermQ{{\textnormal{Q}}}
\def\ermR{{\textnormal{R}}}
\def\ermS{{\textnormal{S}}}
\def\ermT{{\textnormal{T}}}
\def\ermU{{\textnormal{U}}}
\def\ermV{{\textnormal{V}}}
\def\ermW{{\textnormal{W}}}
\def\ermX{{\textnormal{X}}}
\def\ermY{{\textnormal{Y}}}
\def\ermZ{{\textnormal{Z}}}

% Vectors
\def\vzero{{\bm{0}}}
\def\vone{{\bm{1}}}
\def\vmu{{\bm{\mu}}}
\def\vtheta{{\bm{\theta}}}
\def\va{{\bm{a}}}
\def\vb{{\bm{b}}}
\def\vc{{\bm{c}}}
\def\vd{{\bm{d}}}
\def\ve{{\bm{e}}}
\def\vf{{\bm{f}}}
\def\vg{{\bm{g}}}
\def\vh{{\bm{h}}}
\def\vi{{\bm{i}}}
\def\vj{{\bm{j}}}
\def\vk{{\bm{k}}}
\def\vl{{\bm{l}}}
\def\vm{{\bm{m}}}
\def\vn{{\bm{n}}}
\def\vo{{\bm{o}}}
\def\vp{{\bm{p}}}
\def\vq{{\bm{q}}}
\def\vr{{\bm{r}}}
\def\vs{{\bm{s}}}
\def\vt{{\bm{t}}}
\def\vu{{\bm{u}}}
\def\vv{{\bm{v}}}
\def\vw{{\bm{w}}}
\def\vx{{\bm{x}}}
\def\vy{{\bm{y}}}
\def\vz{{\bm{z}}}

% Elements of vectors
\def\evalpha{{\alpha}}
\def\evbeta{{\beta}}
\def\evepsilon{{\epsilon}}
\def\evlambda{{\lambda}}
\def\evomega{{\omega}}
\def\evmu{{\mu}}
\def\evpsi{{\psi}}
\def\evsigma{{\sigma}}
\def\evtheta{{\theta}}
\def\eva{{a}}
\def\evb{{b}}
\def\evc{{c}}
\def\evd{{d}}
\def\eve{{e}}
\def\evf{{f}}
\def\evg{{g}}
\def\evh{{h}}
\def\evi{{i}}
\def\evj{{j}}
\def\evk{{k}}
\def\evl{{l}}
\def\evm{{m}}
\def\evn{{n}}
\def\evo{{o}}
\def\evp{{p}}
\def\evq{{q}}
\def\evr{{r}}
\def\evs{{s}}
\def\evt{{t}}
\def\evu{{u}}
\def\evv{{v}}
\def\evw{{w}}
\def\evx{{x}}
\def\evy{{y}}
\def\evz{{z}}

% Matrix
\def\mA{{\bm{A}}}
\def\mB{{\bm{B}}}
\def\mC{{\bm{C}}}
\def\mD{{\bm{D}}}
\def\mE{{\bm{E}}}
\def\mF{{\bm{F}}}
\def\mG{{\bm{G}}}
\def\mH{{\bm{H}}}
\def\mI{{\bm{I}}}
\def\mJ{{\bm{J}}}
\def\mK{{\bm{K}}}
\def\mL{{\bm{L}}}
\def\mM{{\bm{M}}}
\def\mN{{\bm{N}}}
\def\mO{{\bm{O}}}
\def\mP{{\bm{P}}}
\def\mQ{{\bm{Q}}}
\def\mR{{\bm{R}}}
\def\mS{{\bm{S}}}
\def\mT{{\bm{T}}}
\def\mU{{\bm{U}}}
\def\mV{{\bm{V}}}
\def\mW{{\bm{W}}}
\def\mX{{\bm{X}}}
\def\mY{{\bm{Y}}}
\def\mZ{{\bm{Z}}}
\def\mBeta{{\bm{\beta}}}
\def\mPhi{{\bm{\Phi}}}
\def\mLambda{{\bm{\Lambda}}}
\def\mSigma{{\bm{\Sigma}}}

% Tensor
\DeclareMathAlphabet{\mathsfit}{\encodingdefault}{\sfdefault}{m}{sl}
\SetMathAlphabet{\mathsfit}{bold}{\encodingdefault}{\sfdefault}{bx}{n}
\newcommand{\tens}[1]{\bm{\mathsfit{#1}}}
\def\tA{{\tens{A}}}
\def\tB{{\tens{B}}}
\def\tC{{\tens{C}}}
\def\tD{{\tens{D}}}
\def\tE{{\tens{E}}}
\def\tF{{\tens{F}}}
\def\tG{{\tens{G}}}
\def\tH{{\tens{H}}}
\def\tI{{\tens{I}}}
\def\tJ{{\tens{J}}}
\def\tK{{\tens{K}}}
\def\tL{{\tens{L}}}
\def\tM{{\tens{M}}}
\def\tN{{\tens{N}}}
\def\tO{{\tens{O}}}
\def\tP{{\tens{P}}}
\def\tQ{{\tens{Q}}}
\def\tR{{\tens{R}}}
\def\tS{{\tens{S}}}
\def\tT{{\tens{T}}}
\def\tU{{\tens{U}}}
\def\tV{{\tens{V}}}
\def\tW{{\tens{W}}}
\def\tX{{\tens{X}}}
\def\tY{{\tens{Y}}}
\def\tZ{{\tens{Z}}}


% Graph
\def\gA{{\mathcal{A}}}
\def\gB{{\mathcal{B}}}
\def\gC{{\mathcal{C}}}
\def\gD{{\mathcal{D}}}
\def\gE{{\mathcal{E}}}
\def\gF{{\mathcal{F}}}
\def\gG{{\mathcal{G}}}
\def\gH{{\mathcal{H}}}
\def\gI{{\mathcal{I}}}
\def\gJ{{\mathcal{J}}}
\def\gK{{\mathcal{K}}}
\def\gL{{\mathcal{L}}}
\def\gM{{\mathcal{M}}}
\def\gN{{\mathcal{N}}}
\def\gO{{\mathcal{O}}}
\def\gP{{\mathcal{P}}}
\def\gQ{{\mathcal{Q}}}
\def\gR{{\mathcal{R}}}
\def\gS{{\mathcal{S}}}
\def\gT{{\mathcal{T}}}
\def\gU{{\mathcal{U}}}
\def\gV{{\mathcal{V}}}
\def\gW{{\mathcal{W}}}
\def\gX{{\mathcal{X}}}
\def\gY{{\mathcal{Y}}}
\def\gZ{{\mathcal{Z}}}

% Sets
\def\sA{{\mathbb{A}}}
\def\sB{{\mathbb{B}}}
\def\sC{{\mathbb{C}}}
\def\sD{{\mathbb{D}}}
% Don't use a set called E, because this would be the same as our symbol
% for expectation.
\def\sF{{\mathbb{F}}}
\def\sG{{\mathbb{G}}}
\def\sH{{\mathbb{H}}}
\def\sI{{\mathbb{I}}}
\def\sJ{{\mathbb{J}}}
\def\sK{{\mathbb{K}}}
\def\sL{{\mathbb{L}}}
\def\sM{{\mathbb{M}}}
\def\sN{{\mathbb{N}}}
\def\sO{{\mathbb{O}}}
\def\sP{{\mathbb{P}}}
\def\sQ{{\mathbb{Q}}}
\def\sR{{\mathbb{R}}}
\def\sS{{\mathbb{S}}}
\def\sT{{\mathbb{T}}}
\def\sU{{\mathbb{U}}}
\def\sV{{\mathbb{V}}}
\def\sW{{\mathbb{W}}}
\def\sX{{\mathbb{X}}}
\def\sY{{\mathbb{Y}}}
\def\sZ{{\mathbb{Z}}}

% Entries of a matrix
\def\emLambda{{\Lambda}}
\def\emA{{A}}
\def\emB{{B}}
\def\emC{{C}}
\def\emD{{D}}
\def\emE{{E}}
\def\emF{{F}}
\def\emG{{G}}
\def\emH{{H}}
\def\emI{{I}}
\def\emJ{{J}}
\def\emK{{K}}
\def\emL{{L}}
\def\emM{{M}}
\def\emN{{N}}
\def\emO{{O}}
\def\emP{{P}}
\def\emQ{{Q}}
\def\emR{{R}}
\def\emS{{S}}
\def\emT{{T}}
\def\emU{{U}}
\def\emV{{V}}
\def\emW{{W}}
\def\emX{{X}}
\def\emY{{Y}}
\def\emZ{{Z}}
\def\emSigma{{\Sigma}}

% entries of a tensor
% Same font as tensor, without \bm wrapper
\newcommand{\etens}[1]{\mathsfit{#1}}
\def\etLambda{{\etens{\Lambda}}}
\def\etA{{\etens{A}}}
\def\etB{{\etens{B}}}
\def\etC{{\etens{C}}}
\def\etD{{\etens{D}}}
\def\etE{{\etens{E}}}
\def\etF{{\etens{F}}}
\def\etG{{\etens{G}}}
\def\etH{{\etens{H}}}
\def\etI{{\etens{I}}}
\def\etJ{{\etens{J}}}
\def\etK{{\etens{K}}}
\def\etL{{\etens{L}}}
\def\etM{{\etens{M}}}
\def\etN{{\etens{N}}}
\def\etO{{\etens{O}}}
\def\etP{{\etens{P}}}
\def\etQ{{\etens{Q}}}
\def\etR{{\etens{R}}}
\def\etS{{\etens{S}}}
\def\etT{{\etens{T}}}
\def\etU{{\etens{U}}}
\def\etV{{\etens{V}}}
\def\etW{{\etens{W}}}
\def\etX{{\etens{X}}}
\def\etY{{\etens{Y}}}
\def\etZ{{\etens{Z}}}

% The true underlying data generating distribution
\newcommand{\pdata}{p_{\rm{data}}}
% The empirical distribution defined by the training set
\newcommand{\ptrain}{\hat{p}_{\rm{data}}}
\newcommand{\Ptrain}{\hat{P}_{\rm{data}}}
% The model distribution
\newcommand{\pmodel}{p_{\rm{model}}}
\newcommand{\Pmodel}{P_{\rm{model}}}
\newcommand{\ptildemodel}{\tilde{p}_{\rm{model}}}
% Stochastic autoencoder distributions
\newcommand{\pencode}{p_{\rm{encoder}}}
\newcommand{\pdecode}{p_{\rm{decoder}}}
\newcommand{\precons}{p_{\rm{reconstruct}}}

\newcommand{\laplace}{\mathrm{Laplace}} % Laplace distribution

\newcommand{\E}{\mathbb{E}}
\newcommand{\Ls}{\mathcal{L}}
\newcommand{\R}{\mathbb{R}}
\newcommand{\emp}{\tilde{p}}
\newcommand{\lr}{\alpha}
\newcommand{\reg}{\lambda}
\newcommand{\rect}{\mathrm{rectifier}}
\newcommand{\softmax}{\mathrm{softmax}}
\newcommand{\sigmoid}{\sigma}
\newcommand{\softplus}{\zeta}
\newcommand{\KL}{D_{\mathrm{KL}}}
\newcommand{\Var}{\mathrm{Var}}
\newcommand{\standarderror}{\mathrm{SE}}
\newcommand{\Cov}{\mathrm{Cov}}
% Wolfram Mathworld says $L^2$ is for function spaces and $\ell^2$ is for vectors
% But then they seem to use $L^2$ for vectors throughout the site, and so does
% wikipedia.
\newcommand{\normlzero}{L^0}
\newcommand{\normlone}{L^1}
\newcommand{\normltwo}{L^2}
\newcommand{\normlp}{L^p}
\newcommand{\normmax}{L^\infty}

\newcommand{\parents}{Pa} % See usage in notation.tex. Chosen to match Daphne's book.

\DeclareMathOperator*{\argmax}{arg\,max}
\DeclareMathOperator*{\argmin}{arg\,min}

\DeclareMathOperator{\sign}{sign}
\DeclareMathOperator{\Tr}{Tr}
\let\ab\allowbreak


\usepackage{hyperref}
\usepackage{url}
\usepackage{booktabs}
\usepackage{amsmath,amssymb,amsthm}
\usepackage{graphicx}
\usepackage{xcolor}

% ---- draft machinery: remove before submission ----
\newcommand{\todo}[1]{\textcolor{red}{[TODO: #1]}}
\newcommand{\placeholder}[1]{\textcolor{blue}{#1}}
% ---------------------------------------------------

\newtheorem{proposition}{Proposition}
\newtheorem{definition}{Definition}

\newcommand{\score}{\mathrm{score}}
\newcommand{\ub}{\mathrm{ub}}
\newcommand{\lb}{\mathrm{lb}}

\title{Certified Robustness of LSTM Autoencoders\\for Time-Series Anomaly Detection}

% Anonymized for double-blind review.
\author{Anonymous authors\\
\texttt{anonymous}}

\newcommand{\fix}{\marginpar{FIX}}
\newcommand{\new}{\marginpar{NEW}}

% \iclrfinalcopy  % uncomment for camera-ready (de-anonymizes)
\begin{document}

\maketitle

\begin{abstract}
Reconstruction-based anomaly detectors---most commonly LSTM
autoencoders---guard safety-critical infrastructure such as electric
power grids, where both failure modes of the detector are attacks: an
adversary may perturb sensor streams to \emph{raise} false alarms, or to
\emph{mask} a genuine anomaly below the detection threshold. Despite a
decade of neural-network verification, no existing deterministic
verifier certifies the defining property of these detectors: a bound on
a \emph{quadratic reconstruction-error score} of a \emph{recurrent
autoencoder} over a perturbation set. We close this gap. We formalize
detector-level certificates---certified false-alarm robustness
($\score \le \tau$ over an $\ell_\infty$ ball around a normal window)
and certified masking robustness ($\score \ge \tau$ around an anomalous
window)---and compute them with sound zonotope abstract transformers for
LSTMs, extending the Cert-RNN framework
to encoder--decoder architectures and quadratic output specifications,
with a predicate-alignment correction that we show is necessary for
soundness on multi-layer models. On LSTM-AE detectors for the IEEE-9
power system, our verifier certifies radii $\sim$\placeholder{100}$\times$
larger than interval bound propagation at equal cost, and matches
CROWN-style linear relaxation to within \placeholder{30\%} of its radius
while being \placeholder{600}$\times$ faster per bound query,
scaling to model sizes where the relaxation exceeds its time budget.
From per-window certificates we derive certified-robust
false-positive/true-positive rates, giving---to our knowledge---the
first deterministic end-to-end guarantee for a learned time-series
anomaly detector. We release the verifier, an adapter harness that runs
all baselines under an identical bisection protocol, and a deterministic
LSTM-AE benchmark generator; no recurrent or anomaly-detection benchmark
has appeared in VNN-COMP to date, and we propose ours to fill that gap.
\end{abstract}

\section{Introduction}
\label{sec:intro}

Learned anomaly detectors increasingly monitor cyber-physical systems:
electric grids, water treatment, industrial control. The dominant
recipe is reconstruction-based: train an autoencoder (typically an LSTM
autoencoder for multivariate time series) on normal operation, and flag
a window $x$ as anomalous when its reconstruction score
$\score(x) = \tfrac{1}{N}\lVert \mathrm{AE}(x) - x\rVert_2^2$ exceeds a
threshold $\tau$ calibrated on training data
\citep{malhotra2016lstm,hundman2018detecting,su2019robust}.
In adversarial settings this detector is itself an attack surface, in
\emph{both} directions. An attacker who can inject small perturbations
into sensor telemetry can (i) push normal windows above $\tau$,
flooding operators with false alarms until the detector is ignored or
disabled, or (ii) shave a genuine fault or attack signature below
$\tau$, hiding it (\emph{anomaly masking}). Empirical attacks of both
kinds are well documented, and recent work argues that reconstruction
detectors are structurally brittle \citep{bouman2025unreliable}. For
infrastructure operators the relevant question is not whether the
detector performs well on average, but what can be \emph{guaranteed}
about its output under bounded sensor perturbation.

Answering that question is a verification problem with three
compounding difficulties. First, the model is \emph{recurrent}:
verification must propagate uncertainty through tens of LSTM steps with
sigmoid/tanh gates and bilinear (gate$\times$state) interactions, where
abstraction error compounds multiplicatively---the regime where
ReLU-specialized feedforward verifiers offer no support. Second, the
architecture is an \emph{autoencoder}: the property couples the
network's output to its \emph{input} ($\mathrm{AE}(x)-x$), unlike the
classification margins all existing RNN verifiers target
\citep{ko2019popqorn,ryou2021scalable,jacoby2020rnnverify,du2021certrnn}.
Third, the specification is \emph{quadratic}: certifying
$\score \le \tau$ (or $\ge \tau$) requires sound bounds on a sum of
squares, not a linear functional of the output. To our knowledge no
published deterministic verifier handles this combination; the only
certified defense for time-series anomaly detection to date is
probabilistic, via randomized smoothing under DTW distance
\citep{liu2026dtw}.

\textbf{This paper.} We certify LSTM autoencoder anomaly detectors
deterministically, end to end, and evaluate what such certification
costs at scale.

\textbf{Contributions.}
\begin{itemize}
\item \textbf{Detector-level certificates.} We formalize two dual
  properties over an $\ell_\infty$ ball of radius $\epsilon$: certified
  \emph{false-alarm} robustness ($\sup \score \le \tau$ around normal
  windows) and certified \emph{masking} robustness
  ($\inf \score \ge \tau$ around anomalous windows), together with
  threat models perturbing one frame or all frames of a window.
  Aggregated over windows these yield \emph{certified-robust FPR/TPR}
  at radius $\epsilon$---the detector-level analogue of certified
  accuracy (\S\ref{sec:method}).
\item \textbf{A sound verifier for recurrent autoencoders.} We extend
  zonotope abstract interpretation for LSTMs
  \citep{du2021certrnn} to encoder--decoder architectures with
  latent broadcast and per-step reconstruction heads, and to sound
  upper \emph{and lower} bounds on the quadratic score. We identify a
  predicate-aliasing unsoundness pattern that arises in multi-layer
  recurrent abstractions under positional generator padding, prove a
  Minkowski-correct alignment fixes it, and show it changes certified
  radii on real checkpoints (\S\ref{sec:method}, \S\ref{sec:results}).
  Our implementation is, to our knowledge, the only public one of the
  Cert-RNN transformer family (the original code was never released).
\item \textbf{A controlled comparison harness.} Certified-radius
  comparisons across tools are confounded by differing search loops and
  spec encodings. Our harness drives every sound tool---our verifier,
  interval bound propagation, and auto\_LiRPA's IBP and CROWN modes on
  the unrolled graph \citep{xu2020automatic}---through the
  \emph{same} radius bisection with the \emph{same} quadratic-score
  assembly from componentwise residual bounds, isolating the abstract
  domain as the only variable; PGD falsification bounds every result
  from above (\S\ref{sec:setup}).
\item \textbf{Benchmarks.} Four trained IEEE-9 power-system LSTM-AE
  checkpoints with a MATLAB star-set cross-reference; a
  \emph{deterministic benchmark generator} producing trained detectors
  at any (hidden width, window length, feature count, depth) from a
  seed; and loaders for public server-telemetry data (SMD). VNN-COMP
  2021--2025 contains no recurrent or anomaly-detection benchmark; we
  propose ours (\S\ref{sec:setup}).
\end{itemize}

\textbf{Findings} (\S\ref{sec:results}). Certification of recurrent
autoencoders exhibits a sharp Pareto structure. Interval propagation is
fast but its bounds explode within a few LSTM steps: at practical
bisection granularity it certifies radii $\sim$\placeholder{100}$\times$
smaller than the zonotope domain, or nothing. CROWN-style backward
relaxation is the tightest sound domain per query---certifying radii
\placeholder{$\sim$1.4}$\times$ ours on small models---but three orders
of magnitude slower, exceeding realistic budgets already at moderate
hidden widths \todo{confirm E2 crossover}. The zonotope verifier holds
the practical middle: \placeholder{sub-second} bound queries at IEEE-9
scale, polynomial growth \todo{fit exponent from E2}, and nonvacuous
certificates on every benchmark. On the detector task, we certify that
\placeholder{X\%} of normal windows admit no false alarm and
\placeholder{Y\%} of anomalous windows cannot be masked at
$\epsilon = $ \placeholder{0.01} \todo{E3}.

\section{Related Work}
\label{sec:related}

\textbf{RNN verification.} POPQORN \citep{ko2019popqorn} bounds RNN,
LSTM and GRU outputs with CROWN-style linear relaxations of the coupled
gate nonlinearities; Prover \citep{ryou2021scalable} learns
polyhedral relaxations optimized per network; RnnVerify
\citep{jacoby2020rnnverify} infers invariants reducing vanilla ReLU-RNN
verification to feedforward queries (LSTMs out of scope); Cert-RNN
\citep{du2021certrnn} propagates zonotopes through sound bilinear and
sigmoid/tanh transformers; star-set methods extend exact/overapproximate
reachability to RNNs and LSTM audio classifiers
\citep{tran2023star,pal2023lstmstar}; abstraction-refinement approaches
followed \citep{lin2026refine}. All target
classification-margin properties of recurrent \emph{classifiers}; none
expresses a reconstruction-error property of an autoencoder, and none of
the first four has maintained public code (POPQORN, Prover, RnnVerify
repositories are unmaintained since 2020--21; Cert-RNN's code was never
released). auto\_LiRPA \citep{xu2020automatic} generalizes linear
relaxation to arbitrary computation graphs including unrolled LSTMs and
is actively maintained within the $\alpha,\beta$-CROWN ecosystem
\citep{wang2021beta,shi2024genbab}; it is our strongest runnable
baseline, given the reconstruction residual as an explicit graph output.

\textbf{Certified defenses for time series and anomaly detection.}
Randomized smoothing yields probabilistic certificates and has recently
been adapted to time-series anomaly detection under DTW perturbations
\citep{liu2026dtw}. Our certificates are deterministic and exact with
respect to the $\ell_\infty$ threat model: no sampling failure
probability, and both directions (false alarm and masking) of the
detector are covered. Deterministic certification of \emph{feedforward}
autoencoder reconstruction appears in reachability analyses of
individual case studies, but we know of no general tool or systematic
evaluation for the recurrent case.

\textbf{Anomaly detection in power systems.} LSTM autoencoders are a
standard detector for PMU/bus telemetry
\citep{malhotra2016lstm,hundman2018detecting}; public power-system
attack datasets exist \citep{adhikari2018power}, and IEEE-bus
simulation is the community-standard data source given CEII
restrictions on utility data. Verification-oriented benchmarks are
absent: no VNN-COMP edition (2021--2025) includes recurrent or
anomaly-detection tasks.

\section{Background}
\label{sec:background}

\textbf{LSTM autoencoder detectors.} An LSTM-AE consists of an encoder
LSTM stack that consumes a window $x = (x_1,\dots,x_T)$,
$x_t \in \mathbb{R}^D$, an (optional multi-layer) decoder LSTM stack
that reads the encoder's final top-layer hidden state $z$ at every step,
and a linear head mapping decoder states back to input space, producing
$\hat{x} = \mathrm{AE}(x)$. The anomaly score is
$\score(x) = \tfrac{1}{N}\lVert \mathrm{AE}(x)-x \rVert_2^2$ with
$N = TD$; the detector flags $x$ iff $\score(x) > \tau$.

\textbf{Zonotopes.} A zonotope
$Z = \{ c + V\alpha : \alpha \in [-1,1]^p \} \subset \mathbb{R}^n$
is closed under affine maps and Minkowski sums, and supports sound
abstract transformers for the LSTM's sigmoid/tanh and bilinear
operations \citep{du2021certrnn}: each transformer returns a zonotope
containing the true image, introducing fresh generators to soundly
overapproximate nonlinear residuals. Crucially, zonotopes are
\emph{relational}: shared generators track first-order correlations
between the perturbed input, the recurrent state, and (for
autoencoders) the subtracted input---exactly what pure interval
propagation discards.

\textbf{Certified radius by bisection.} Given a sound predicate
``property holds over the $\epsilon$-ball,'' Algorithm~1 of
\citet{du2021certrnn} performs a $k$-step bisection on $\epsilon$
returning the largest certified radius at granularity
$\epsilon_0 2^{-k}$. All tools in our study answer the same predicate
inside the same bisection (\S\ref{sec:setup}).

\section{Certifying reconstruction-error properties}
\label{sec:method}

\subsection{Properties and threat models}

\begin{definition}[False-alarm robustness]
A detector $(\mathrm{AE},\tau)$ is $\epsilon$-certified
\emph{false-alarm robust} at a normal window $x^\ast$ iff
$\sup_{\lVert x' - x^\ast \rVert_\infty \le \epsilon} \score(x') \le \tau$.
\end{definition}

\begin{definition}[Masking robustness]
A detector is $\epsilon$-certified \emph{masking robust} at an
anomalous window $x^\ast$ iff
$\inf_{\lVert x' - x^\ast\rVert_\infty \le \epsilon} \score(x') \ge \tau$.
\end{definition}

We instantiate the ball under two threat models: \emph{single-frame}
(one time step perturbed; per-frame radii reported as the minimum over
frames) and \emph{multi-frame} (all steps perturbed jointly---the
operationally conservative model for an attacker with stream access).
Over a set of windows, the fraction whose certified radius covers a
target $\epsilon$ gives \emph{certified-robust FPR} (normal windows)
and \emph{certified-robust TPR} (anomalous windows)---the direct
analogue of certified accuracy for classifiers, at detector level.

\subsection{Abstract autoencoder forward pass}

We propagate the input zonotope through the encoder stack, broadcast
the final top-layer hidden zonotope $Z_z$ into every decoder step, and
apply the head per step, yielding reconstruction zonotopes
$Z_{\hat{x}_t}$. Because the input frames are themselves zonotopes
sharing generators with the network state, the residual
$Z_{r_t} = Z_{\hat{x}_t} \ominus Z_{x_t}$ is computed \emph{in the
shared generator space}, retaining input--output correlation; this is
the step interval arithmetic cannot express and a key source of our
advantage over IBP (\S\ref{sec:results}).

\textbf{Sound score bounds.} With componentwise residual ranges
$[l_{td}, u_{td}]$ from $Z_{r_t}$,
\begin{equation}
\label{eq:bounds}
\score_{\ub} = \tfrac{1}{N}\textstyle\sum_{t,d} \max(|l_{td}|,|u_{td}|)^2,
\qquad
\score_{\lb} = \tfrac{1}{N}\textstyle\sum_{t,d} \max(0, l_{td}, -u_{td})^2,
\end{equation}
which are sound for the sup/inf of the score since a sum's extremum is
bounded by the sum of componentwise extrema. $\score_{\ub} \le \tau$
certifies false-alarm robustness; $\score_{\lb} \ge \tau$ certifies
masking robustness. The looseness of \eqref{eq:bounds} relative to a
joint quadratic optimization over the zonotope is measured empirically
in \S\ref{sec:results} \todo{E4 gap numbers}; tightening it (e.g.\ by
quadratic programming over the generator box) is a natural next step.

\subsection{Predicate alignment, and an unsoundness pattern in the wild}
\label{sec:aliasing}

Zonotope transformers for recurrent networks repeatedly combine
zonotopes born at different times: the previous state (carrying fresh
generators from step $t{-}1$) and the current gate pre-activations
(carrying fresh generators from step $t$). If implementations align
generator matrices \emph{positionally} (zero-padding to equal width),
two distinct noise symbols can silently share a column---and by the
triangle inequality the merged $|V|$-sums \emph{shrink}, producing
bounds tighter than the true reachable set: an unsound certificate.
This is not hypothetical: the MATLAB star-set reference implementation
we cross-validate against exhibits exactly this pattern on multi-layer
models, certifying radii up to \placeholder{7\%} larger than the sound
value on every frame of our two-layer IEEE-9 checkpoint
(\S\ref{sec:results}). Our implementation assigns every generator a
persistent identity and embeds operands into a unified predicate space
(shared ids share columns; unshared get disjoint columns), which we
validate by exhaustive per-transformer fuzzing against LP ground truth
and end-to-end PGD probing. \todo{cite soundness proofs in appendix;
port docs/soundness.md}

\subsection{Complexity}

For hidden width $H$, window length $T$, depth $L$: each LSTM step
applies $O(H)$ bilinear transformers, each touching generator columns
accumulated so far; fresh generators grow by $O(H)$ per (step, layer),
so the abstract forward pass costs
$O\!\big(T^2 L^2 H^2 (D+H)\big)$ arithmetic in the worst case
\todo{verify constants; measure empirical exponents in E2}, against
$O(TLH(D+H))$ for one concrete forward. The bisection multiplies either
by $k{+}1$ (multi-frame) or $T(k{+}1)$ (single-frame, all frames).

\section{Experimental setup}
\label{sec:setup}

\textbf{Tools.} (i) \textsc{Zono} (ours); (ii) \textsc{Ibp}: interval
bound propagation implemented on the same model dicts (endpoint
propagation through monotone gates, four-product interval
multiplication); (iii) \textsc{Lirpa-Ibp} and (iv) \textsc{Crown}:
auto\_LiRPA \citep{xu2020automatic} interval and backward modes on the
explicit-op unrolled AE graph (fp32, its backward-mode requirement;
all other tools fp64); (v) \textsc{Pgd}: multi-restart projected
gradient attack on the concrete score, giving an empirical upper bound
on every radius. Table~\ref{tab:tools} situates the non-runnable RNN
verifiers: none supports the reconstruction-error property, and none
has maintained code.

\begin{table}[h]
\caption{RNN verification tools vs.\ the requirements of the LSTM-AE
detector property (as of 2026-08).}
\label{tab:tools}
\begin{center}
\small
\begin{tabular}{lcccc}
\toprule
 & LSTM & AE recon.\ spec & quadratic score & public code maintained \\
\midrule
POPQORN \citep{ko2019popqorn}      & \checkmark & -- & -- & dead (2020) \\
Prover \citep{ryou2021scalable}    & \checkmark & -- & -- & dead (2021) \\
RnnVerify \citep{jacoby2020rnnverify} & -- & -- & -- & dead (2020) \\
Cert-RNN \citep{du2021certrnn}     & \checkmark & -- & -- & never released \\
auto\_LiRPA \citep{xu2020automatic} & \checkmark & graph-expressible & assembled \eqref{eq:bounds} & active \\
\textbf{Ours}                      & \checkmark & \checkmark & \checkmark\,(ub+lb) & released \\
\bottomrule
\end{tabular}
\end{center}
\end{table}

\textbf{Fairness protocol.} Every sound tool answers the same predicate
inside the same bisection ($\epsilon_0 = 0.5$, $k=12$ unless noted),
and the quadratic score is assembled from componentwise residual
bounds identically \eqref{eq:bounds} for all tools---so measured
differences isolate the abstract domain. Single-thread BLAS; identical
hardware (28-core Linux workstation); wall-clock and bound-call counts
recorded per cell; JSONL results with environment stamps are shipped
with the code.

\textbf{Benchmarks.} (a) \emph{IEEE-9}: four LSTM-AE checkpoints
trained on normalized IEEE-9 bus telemetry ($D{=}36$, $T{=}30$):
S/M/L single-layer with $H \in \{4,16,64\}$ and D two-layer with
$H{=}16$; thresholds $\tau$ calibrated on training scores; a MATLAB
star-set implementation provides an independent reference (and the
aliasing case study). (b) \emph{Synthetic lattice}: a deterministic
generator trains detectors on seeded quasi-periodic multivariate
processes at any $(H,T,D,L)$; sweeps vary one axis around center
$(16,30,9,1)$: $H \le 128$, $T \le 120$, $D \le 72$, $L \le 3$; injected
spike/step/dropout anomalies (flagged by the clean detector) provide
masking subjects. (c) \emph{SMD} \citep{su2019robust}: server telemetry
($D{=}38$), standard TSAD benchmark, for cross-domain generality.
\todo{wire Morris power-attack PMU dataset \citep{adhikari2018power}}

\section{Results}
\label{sec:results}

\subsection{E1: Verifiability across tools}
\todo{full table from results/e1.jsonl; smoke row below}

\begin{table}[h]
\caption{Certified radius (single-frame, min over frames) and cost.
Preliminary smoke cell, IEEE-9~S frame~0, $k{=}8$: \textsc{Zono}
0.0156 at 0.05\,s/query; \textsc{Crown} 0.0215 at $\sim$28\,s/query;
\textsc{Ibp}/\textsc{Lirpa-Ibp} certify nothing at $2^{-8}$
granularity; \textsc{Pgd} finds no violation up to $\epsilon \approx 1$.}
\label{tab:e1}
\begin{center}
\begin{tabular}{lccccc}
\toprule
 & \textsc{Zono} & \textsc{Ibp} & \textsc{Lirpa-Ibp} & \textsc{Crown} & \textsc{Pgd} (ub) \\
\midrule
IEEE9-S & \todo{} & & & & \\
IEEE9-M & & & & & \\
IEEE9-L & & & & & \\
IEEE9-D & & & & & \\
SMD     & & & & & \\
\bottomrule
\end{tabular}
\end{center}
\end{table}

\subsection{E2: Scalability}
\todo{scaling figures (sec/query vs H, T, D, L, log-log) +
cactus plot; report empirical exponents and CROWN's budget crossover.}

\subsection{E3: Certified detection}
\todo{certified-robust FPR/TPR at target radii; both properties;
multi-frame threat model; synthetic + SMD.}

\subsection{E4: Tightness and ablations}
\todo{(a) certified-vs-PGD gap decomposition (domain looseness vs
quadratic-bound looseness); (b) radius/time vs bisection depth;
(c) single- vs multi-frame radii; (d) the aliasing case study: MATLAB
reference certifies $+1.7$--$2.7\times 10^{-3}$ larger $\epsilon$ on
all 30 frames of IEEE9-D --- theoretically unsound bounds that PGD
cannot falsify because of quadratic-bound slack; discussion.}

\section{Limitations}
\label{sec:limitations}

Our certificates are $\ell_\infty$-ball guarantees around individual
windows: they do not cover distribution shift, semantic attacks outside
the ball, or windows not analyzed. The componentwise quadratic bound
\eqref{eq:bounds} leaves measurable slack (E4); the auto\_LiRPA
comparison runs in fp32 by that library's requirement; and our
synthetic lattice, while controlled and reproducible, is a proxy---the
public-data results carry the external-validity weight. Masking
certificates hold at the analyzed anomalous windows and do not preclude
attacks that suppress anomalies \emph{before} they manifest in the
window.

\section{Conclusion}

We gave the first deterministic certification pipeline for LSTM
autoencoder anomaly detectors---both false-alarm and masking
robustness---together with a controlled cross-tool evaluation and
reproducible benchmarks. The zonotope domain occupies the practical
sweet spot between interval collapse and linear-relaxation cost, and
detector-level certified FPR/TPR turns per-window radii into an
operator-meaningful guarantee. We propose the benchmark suite for
community adoption (e.g.\ VNN-COMP, which to date includes no recurrent
or anomaly-detection tasks).

\subsubsection*{Reproducibility statement}
All experiments are driven by a results-record harness (JSONL with
hardware/software stamps) shipped with the code; the synthetic
benchmark family regenerates byte-identically from seeds; IEEE-9
checkpoints, anchors, and thresholds are released. \todo{anonymized
code link}

\bibliography{refs}
\bibliographystyle{iclr2027_conference}

\appendix
\section{Soundness of the transformers}
\todo{port the transformer soundness statements and proofs
(bilinear sigmoid-identity / sigmoid-tanh, alignment lemma) from
docs/soundness.md.}

\section{The predicate-aliasing pattern}
\todo{port the mechanism (docs/lstm\_ae\_results.md \S\,``Why D
diverges''): distinct fresh generators sharing a column index under
positional padding; triangle-inequality shrinkage; per-frame deltas.}

\section{Experiment details}
\todo{hyperparameters, benchmark generator process, PGD config,
timeout policy, hardware.}

\end{document}
